\documentclass[11pt,letterpaper]{article}
\usepackage[utf8]{inputenc}
\usepackage[T1]{fontenc}
\usepackage[spanish,es-nodecimaldot,es-tabla]{babel}
\usepackage{mathptmx}
\usepackage{amsmath,amssymb}
\usepackage{array,booktabs}
\usepackage{xcolor}
\usepackage[landscape,letterpaper,left=1.3cm,right=1.3cm,top=0.9cm,bottom=0.9cm]{geometry}
\setlength{\parindent}{0pt}\setlength{\parskip}{0.3em}
\newcommand{\chk}{\framebox[0.42cm]{\rule{0pt}{0.30cm}}}
\pagestyle{empty}
\begin{document}
\begin{center}
{\large \textbf{Pauta Cátedra --- Caso 1: Regresión}}\\[0.1cm]
{\footnotesize Proyecto Integrador · Estadística Aplicada · Marque el nivel logrado en cada dimensión (uso del profesor)}
\end{center}
\vspace{0.05cm}
{\small Equipo: \underline{\hspace{5cm}}\quad Proyecto: \underline{\hspace{5.5cm}}\quad Fecha: \underline{\hspace{2.6cm}}\quad Evaluador: \underline{\hspace{3.5cm}}}
\vspace{0.15cm}

\footnotesize\renewcommand{\arraystretch}{1.15}
\begin{tabular}{p{2.4cm} p{3.15cm} p{3.15cm} p{3.15cm} p{3.15cm} p{3.15cm} p{1.2cm}}
\toprule
Dimensión & Nivel 1 (20\%) & Nivel 2 (40\%) & Nivel 3 (60\%) & Nivel 4 (80\%) & Nivel 5 (100\%) & Obt. \\
\midrule
Oratoria y dominio conceptual (45) & \chk\ \textbf{9}~~Lee la mayor parte de la exposición o no logra explicar; no responde las preguntas o comete errores conceptuales graves. & \chk\ \textbf{18}~~Recita apoyándose en notas o láminas; responde de forma vaga o imprecisa; dominio conceptual parcial. & \chk\ \textbf{27}~~Expone con apoyo puntual; explica las ideas centrales y responde lo básico, con vacíos menores. & \chk\ \textbf{36}~~Expone con fluidez sin leer, explica en vez de recitar y responde con solvencia la mayoría de las preguntas. & \chk\ \textbf{45}~~Domina y conecta los conceptos con la evidencia y responde con precisión preguntas fuera de lo previsto, de forma individual. & \rule{0pt}{0.7cm} \\
Calidad de la presentación (35) & \chk\ \textbf{7}~~Láminas sobrecargadas o ilegibles, con volcados de código o texto; sin narrativa. & \chk\ \textbf{14}~~Diseño irregular, varias ideas por lámina o figuras sin título ni comentario; narrativa débil. & \chk\ \textbf{21}~~Diseño funcional, una idea por lámina en general y figuras legibles; narrativa comprensible. & \chk\ \textbf{28}~~Diseño claro y consistente, una idea por lámina, figuras con título, unidades y comentario; narrativa coherente. & \chk\ \textbf{35}~~Diseño sobresaliente que refuerza el mensaje; cada lámina aporta una conclusión visible y la historia se sigue sin esfuerzo. & \rule{0pt}{0.7cm} \\
Rigor técnico y decisiones (15) & \chk\ \textbf{3}~~Método incorrecto o sin baseline; conclusiones no sostenibles o causalidad injustificada. & \chk\ \textbf{6}~~Método incompleto; comparaciones sin incertidumbre o decisiones poco justificadas. & \chk\ \textbf{9}~~Método correcto en lo esencial (validación, baseline, métricas) con justificación limitada. & \chk\ \textbf{12}~~Decisiones justificadas y coherentes con lo presentado; métricas e interpretación adecuadas. & \chk\ \textbf{15}~~Además, discute supuestos, sensibilidad y límites, y descarta complejidad innecesaria con evidencia. & \rule{0pt}{0.7cm} \\
Informe de respaldo (5) & \chk\ \textbf{1}~~Ausente, excede el límite de páginas o contradice la presentación o el script. & \chk\ \textbf{2}~~Incompleto o meramente descriptivo; respalda débilmente las afirmaciones. & \chk\ \textbf{3}~~Cumple el formato ($\leq 5$ páginas) y respalda las afirmaciones principales. & \chk\ \textbf{4}~~Claro, trazable y coherente con la presentación y el script; tablas y figuras sustentan decisiones. & \chk\ \textbf{5}~~Síntesis precisa que permite reconstruir la evidencia y la recomendación. & \rule{0pt}{0.7cm} \\
\bottomrule
\end{tabular}

\vspace{0.2cm}
{\small
En cada celda, el número en negrita es el puntaje de esa dimensión en ese nivel (Nivel $k = k\times 20\%$ de logro). \textbf{Subtotal} (suma de las casillas marcadas, máx. 100): \underline{\hspace{2cm}}\\[0.2cm]
\textbf{Penalizaciones (descuentos fijos, acumulables):}\quad \chk\ Lee gran parte de la exposición (láminas o notas) $(-15)$\quad \chk\ No maneja los conceptos al ser preguntado $(-15)$\quad \chk\ Informe ausente o supera 5 páginas $(-5)$\quad \chk\ Excede el tiempo (>10 min) o no participan todos los integrantes $(-5)$\quad \textbf{Descuento:} \underline{\hspace{1.6cm}}\\[0.2cm]
\textbf{Co-evaluación:}\quad \% de logro recibido: \underline{\hspace{1.8cm}}\quad (si $<60\%$: nota individual $=$ nota del grupo $\times$ logro; sin respuesta $=$ nota mínima). Aplica a ambas notas.\\[0.2cm]
\textbf{NOTA CÁTEDRA} (subtotal $-$ descuentos): \underline{\hspace{2.2cm}}
}

\clearpage
\begin{center}
{\large \textbf{Pauta Laboratorio --- Caso 1: Regresión}}\\[0.1cm]
{\footnotesize Proyecto Integrador · Estadística Aplicada · Marque el nivel logrado en cada dimensión (uso del profesor)}
\end{center}
\vspace{0.05cm}
{\small Equipo: \underline{\hspace{5cm}}\quad Proyecto: \underline{\hspace{5.5cm}}\quad Fecha: \underline{\hspace{2.6cm}}\quad Evaluador: \underline{\hspace{3.5cm}}}
\vspace{0.15cm}

\footnotesize\renewcommand{\arraystretch}{1.15}
\begin{tabular}{p{2.4cm} p{3.15cm} p{3.15cm} p{3.15cm} p{3.15cm} p{3.15cm} p{1.2cm}}
\toprule
Dimensión & Nivel 1 (20\%) & Nivel 2 (40\%) & Nivel 3 (60\%) & Nivel 4 (80\%) & Nivel 5 (100\%) & Obt. \\
\midrule
Reproducibilidad y ejecución (40) & \chk\ \textbf{8}~~No ejecuta o no reproduce lo presentado; requiere reescritura. & \chk\ \textbf{16}~~Ejecuta con intervención manual no documentada o reproduce sólo parcialmente. & \chk\ \textbf{24}~~Reproduce los resultados principales desde una sesión limpia siguiendo instrucciones básicas. & \chk\ \textbf{32}~~Reproduce todo con una ejecución única y documentada. & \chk\ \textbf{40}~~Además, automatiza la ejecución y verifica la reproducibilidad (semillas globales, comprobaciones). & \rule{0pt}{0.7cm} \\
Pipeline sin fuga (leakage) (35) & \chk\ \textbf{7}~~Fuga evidente: preprocesamiento global, o el test participa del ajuste o la selección. & \chk\ \textbf{14}~~Validación con errores u operaciones fuera de los folds que alteran la estimación. & \chk\ \textbf{21}~~Separación entrenamiento-validación-test correcta y validación cruzada cuando corresponde. & \chk\ \textbf{28}~~Toda transformación, selección y remuestreo ocurren dentro de cada fold; el test permanece intacto. & \chk\ \textbf{35}~~Además, incluye verificaciones que detectan fuga, inconsistencias de partición o fuga temporal. & \rule{0pt}{0.7cm} \\
Legibilidad, semillas y coherencia (25) & \chk\ \textbf{5}~~Código ilegible o sin semillas; no corresponde con lo presentado. & \chk\ \textbf{10}~~Nombres inconsistentes o documentación mínima; resultados que no coinciden con la presentación. & \chk\ \textbf{15}~~Código legible, con semillas y documentación básica; coincide en lo esencial con la presentación. & \chk\ \textbf{20}~~Modular y bien documentado; los resultados coinciden exactamente con la presentación y el informe. & \chk\ \textbf{25}~~Además, es parametrizable, con manejo de errores y registro de configuraciones. & \rule{0pt}{0.7cm} \\
\bottomrule
\end{tabular}

\vspace{0.2cm}
{\small
En cada celda, el número en negrita es el puntaje de esa dimensión en ese nivel (Nivel $k = k\times 20\%$ de logro). \textbf{Subtotal} (suma de las casillas marcadas, máx. 100): \underline{\hspace{2cm}}\\[0.2cm]
\textbf{Penalizaciones (descuentos fijos, acumulables):}\quad \chk\ Fuga de información (leakage) o uso del test para seleccionar $(-25)$\quad \chk\ Requiere cambios manuales no documentados para ejecutar $(-10)$\quad \textbf{Descuento:} \underline{\hspace{1.6cm}}\\[0.2cm]
\textbf{NOTA LABORATORIO} (subtotal $-$ descuentos): \underline{\hspace{2.2cm}}
}

\end{document}
